\documentclass[12pt]{article}
\usepackage[utf8]{inputenc}
\usepackage[T1]{fontenc}
\usepackage{hyperref}
\title{Tsum Language Resources for K-12 Curriculum Development}
\author{Compiled by Claude Sonnet 5 (claude-sonnet-5) for Andr\'as Kornai}
\date{}
\begin{document}
\maketitle

\section{Introduction}
Tsum (ISO 639-3: ttz; also Tsumke) is a Tibetic language of Gorkha District, Nepal, with about 4,800 speakers (2000 census, as cited by Wikipedia). It is a two-tone language. Younger speakers are increasingly educated in Nepali and English, and the language is attriting. Search results showed little material designed for schools.

\section{Primary Educational Resources}
No primers, textbooks or school materials for Tsum were found in the search results.

\section{Linguistic and Reference Materials}
\begin{itemize}
\item Overview of the language: \url{https://en.wikipedia.org/wiki/Tsum_language}
\item Donohue, M. and Dhakal, D. N., \emph{A Tsum Lexicon} (LINCOM, 2016, 197 pp.). This is the main lexical resource; it was identified in search results, but no direct link to the book was seen.
\item Naijing Liu, MA thesis \emph{Tsum tone: a challenge for typology and phonological description} (ANU, 2015); mentioned in search results, no direct URL seen.
\item ASJP wordlist: \url{https://asjp.clld.org/languages/TSUM}
\end{itemize}
No full published grammar was found.

\section{Religious and Community Resources}
None were found in the search results.

\section{Digital and Audio Resources}
Naijing Liu's Tsum Data Collection (NL1) at PARADISEC holds audio and video recordings (verb morphosyntax, nominal morphology, reduplication). Example items:
\begin{itemize}
\item \url{https://catalog.paradisec.org.au/collections/NL1/items/ttz0029}
\item \url{https://catalog.paradisec.org.au/collections/NL1/items/ttz0069}
\item \url{https://catalog.paradisec.org.au/collections/NL1/items/ttz0077}
\end{itemize}

\section{Curriculum Development Context}
Documentation consists of a lexicon, a tone study and archived recordings. A curriculum would have to be built largely from scratch, starting from the lexicon and the archived recordings. A sociolinguistic survey by Webster (2022, \emph{Journal of Language Survey Reports}) covers North Gorkha dialects including Tsum; no URL was seen.

\end{document}